\subsection{连接同态与正合序列}

设M是一个右A-模。回忆对于每个左A-模N，我们在前一小节（X，第69页，命题5）中定义了一个同构

设

\[
\psi_ {\mathbf {M}} (\mathbf {N}): \operatorname{Tor} ^ {\mathbf {A}} (\mathbf {M}, \mathbf {N}) \rightarrow \mathrm{H} (\mathrm{L} (\mathbf {M}) \otimes_ {\mathbf {A}} \mathbf {N}).\tag{ε}
\]

\[
0 \rightarrow \mathbf {N} ^ {\prime} \xrightarrow {u} \mathbf {N} \xrightarrow {v} \mathbf {N} ^ {\prime \prime} \rightarrow 0
\]

一个左A-模的正合序列；则 \(k\) -模复形的序列

\[
\left(^ {\mathrm{M}} \mathcal {E}\right) \quad 0 \longrightarrow \mathrm{L} (\mathrm{M}) \otimes_ {\mathrm{A}} \mathrm{N} ^ {\prime} \xrightarrow {1 \otimes u} \mathrm{L} (\mathrm{M}) \otimes_ {\mathrm{A}} \mathrm{N} \xrightarrow {1 \otimes v} \mathrm{L} (\mathrm{M}) \otimes_ {\mathrm{A}} \mathrm{N} ^ {\prime \prime} \longrightarrow 0
\]

这一序列是正合的（X，第66页，引理1）；令

\[
\hat {c} \left(^ {\mathrm{M}} \mathcal {E}\right): \mathrm{H} (\mathrm{L} (\mathrm{M}) \otimes_ {\mathrm{A}} \mathrm{N} ^ {\prime \prime}) \rightarrow \mathrm{H} (\mathrm{L} (\mathrm{M}) \otimes_ {\mathrm{A}} \mathrm{N} ^ {\prime})
\]

为相应的连接同态（X，第 29 页）。

定义 2. --- 相对于模 M 和正合序列 \(\mathcal{E}\) 的挠积的连接同态指复合同态

\[
\partial (\mathrm{M}, \mathcal {E}) = \psi_ {\mathrm{M}} \left(\mathrm{N} ^ {\prime}\right) ^ {- 1} \circ \partial \left(^ {\mathrm{M}} \mathcal {E}\right) \circ \psi_ {\mathrm{M}} \left(\mathrm{N} ^ {\prime \prime}\right): \operatorname{Tor} ^ {\mathrm{A}} \left(\mathrm{M}, \mathrm{N} ^ {\prime \prime}\right)\rightarrow \operatorname{Tor} ^ {\mathrm{A}} \left(\mathrm{M}, \mathrm{N} ^ {\prime}\right).
\]

这是一个分次 \(k\) -同态，次数为 \((-1)\) ，其齐次分量记为 \(\partial_n(\mathbf{M}, \mathcal{E}) : \mathrm{Tor}_n^{\mathbf{A}}(\mathbf{M}, \mathbf{N}'') \to \mathrm{Tor}_{n-1}^{\mathbf{A}}(\mathbf{M}, \mathbf{N}')\) .

定理 1. --- k-模同态的向左无界序列

\[
\begin{array}{l} \dots \longrightarrow \operatorname{Tor} _ {n} ^ {\mathbf {A}} (\mathbf {M}, \mathbf {N} ^ {\prime}) \xrightarrow {\operatorname{Tor} _ {n} ^ {\mathbf {A}} (1 , u)} \operatorname{Tor} _ {n} ^ {\mathbf {A}} (\mathbf {M}, \mathbf {N}) \xrightarrow {\operatorname{Tor} _ {n} ^ {\mathbf {A}} (1 , v)} \operatorname{Tor} _ {n} ^ {\mathbf {A}} (\mathbf {M}, \mathbf {N} ^ {\prime \prime}) \\ \xrightarrow {\hat {c} _ {n} (\mathbf {M} , \mathcal {E})} \operatorname{Tor} _ {n - 1} ^ {\mathbf {A}} (\mathbf {M}, \mathbf {N} ^ {\prime}) \xrightarrow {\operatorname{Tor} _ {n - 1} ^ {\mathbf {A}} (1 , u)} \dots \xrightarrow {\operatorname{Tor} _ {1} ^ {\mathbf {A}} (1 , v)} \operatorname{Tor} _ {1} ^ {\mathbf {A}} (\mathbf {M}, \mathbf {N} ^ {\prime \prime}) \\ \xrightarrow {\hat {c} _ {1} (\mathbf {M} , \mathcal {E})} \mathbf {M} \otimes_ {\mathbf {A}} \mathbf {N} ^ {\prime} \xrightarrow {1 \otimes u} \mathbf {M} \otimes_ {\mathbf {A}} \mathbf {N} \xrightarrow {1 \otimes v} \mathbf {M} \otimes_ {\mathbf {A}} \mathbf {N} ^ {\prime \prime} \longrightarrow 0 \end{array}
\]

是正合的。

事实上，考虑图表

\[
\begin{array}{c} \operatorname{Tor} (\mathbf {M}, \mathbf {N} ^ {\prime}) \xrightarrow {\operatorname{Tor} (1 , u)} \operatorname{Tor} (\mathbf {M}, \mathbf {N}) \xrightarrow {\operatorname{Tor} (1 , v)} \operatorname{Tor} (\mathbf {M}, \mathbf {N} ^ {\prime \prime}) \xrightarrow {\partial (\mathbf {M} , \mathcal {E})} \operatorname{Tor} (\mathbf {M}, \mathbf {N} ^ {\prime}) \xrightarrow {\operatorname{Tor} (1 , u)} \operatorname{Tor} (\mathbf {M}, \mathbf {N}) \\ \psi_ {\mathbf {M} (\mathbf {N} ^ {\prime})} \Big | _ {\downarrow} \quad \psi_ {\mathbf {M} (\mathbf {N})} \Big | _ {\downarrow} \quad \psi_ {\mathbf {M} (\mathbf {N} ^ {\prime \prime})} \Big | _ {\downarrow} \quad \psi_ {\mathbf {M} (\mathbf {N} ^ {\prime})} \Big | _ {\downarrow} \quad \psi_ {\mathbf {M} (\mathbf {N} ^ {\prime})} \Big | _ {\downarrow} \\ H (L (\mathbf {M}) \otimes N ^ {\prime}) \xrightarrow {H (1 \otimes u)} H (L (\mathbf {M}) \otimes N) \xrightarrow {H (1 \otimes v)} H (L (\mathbf {M} \otimes N ^ {\prime \prime}) \xrightarrow {\partial^ {(M)} \mathcal {E}} H (L (\mathbf {M}) \otimes N ^ {\prime}) \xrightarrow {} H (L (\mathbf {M}) \otimes N). \end{array}
\]

由（X，第 69 页，注记 3）和定义 2 知它是交换的。另一方面，底行是正合的（X，第 30 页，定理 1），并且各个 \(\psi_{\mathrm{M}}\) 是双射的（X，第 69 页，命题 5）。

推论 1. --- 如果 Tor \(_{1}^{A}\) (M, N'') = 0，则序列

\[
0 \longrightarrow \mathrm{M} \otimes_ {\mathrm{A}} \mathrm{N} ^ {\prime} \xrightarrow {1 \otimes u} \mathrm{M} \otimes_ {\mathrm{A}} \mathrm{N} \xrightarrow {1 \otimes v} \mathrm{M} \otimes_ {\mathrm{A}} \mathrm{N} ^ {\prime \prime} \longrightarrow 0
\]

是正合的。

推论 2. --- 设 \(0 \to C' \xrightarrow{u} C \xrightarrow{v} C'' \to 0\) 是左 A-模复形的正合序列，且 E 是右 A-模复形。如果 \(C''\) 或 E 是平坦的，则序列

\[
0 \longrightarrow \mathrm{E} \otimes_ {\mathrm{A}} \mathrm{C} ^ {\prime} \xrightarrow {1 \otimes u} \mathrm{E} \otimes_ {\mathrm{A}} \mathrm{C} \xrightarrow {1 \otimes v} \mathrm{E} \otimes_ {\mathrm{A}} \mathrm{C} ^ {\prime \prime} \longrightarrow 0
\]

是正合的。

事实上， \(\operatorname{Tor}_1^{\mathbf{A}}(\mathbf{E},\mathbf{C}^{\prime \prime}) = 0\) 根据 X，第 69 页，命题 5 的推论。

例. --- 设 a 是 A 的一个理想。左 A-模的正合序列

\[
0 \rightarrow \mathfrak {a} \rightarrow \mathrm{A} _ {\mathrm{s}} \rightarrow \mathrm{A} / \mathfrak {a} \rightarrow 0
\]

给出挠积的正合序列，其中项 \(\operatorname{Tor}_{i}^{\mathbf{A}}(\mathbf{M},\mathbf{A})\) 对于 i \textgreater{} 0 为零。由此我们推出同构

\[
\operatorname{Tor} _ {i + 1} ^ {\mathbf {A}} (\mathbf {M}, \mathbf {A} / \mathfrak {a}) \rightarrow \operatorname{Tor} _ {i} ^ {\mathbf {A}} (\mathbf {M}, \mathfrak {a}), \quad i > 0
\]

以及一个正合序列

\[
0 \rightarrow \operatorname{Tor} _ {1} ^ {\mathbf {A}} (\mathrm{M}, \mathrm{A} / \mathfrak {a}) \rightarrow \mathrm{M} \otimes_ {\mathbf {A}} \mathfrak {a} \rightarrow \mathrm{M} \otimes_ {\mathbf {A}} \mathrm{A} \rightarrow \mathrm{M} \otimes \mathrm{A} / \mathfrak {a} \rightarrow 0:
\]

由此可知 \(\operatorname{Tor}_{1}^{\mathbf{A}}(\mathbf{M},\mathbf{A}/\mathfrak{a})\) 与典范同态 \(\mathbf{M}\otimes_{\mathbf{A}}\mathfrak{a}\to\mathbf{M}\) 的核等同。

例如，取M为如下形式的模： \(A_{d}/b\) ，其中 b 是 A 的一个右理想，则得到 \(\operatorname{Tor}_{1}^{A}(A/b, A/a)\) 到 \((a \cap b)/ba\) 上的一个同构。

命题 9. --- 设 \(f: \mathbf{M} \to \mathbf{M}_{1}\) 是右A-模的同态，且

(ε)

\[
\begin{array}{c} 0 \longrightarrow \mathrm{N} ^ {\prime} \stackrel {{u}} {{\longrightarrow}} \mathrm{N} \stackrel {{v}} {{\longrightarrow}} \mathrm{N} ^ {\prime \prime} \longrightarrow 0 \\ g ^ {\prime} \Biggl \downarrow \qquad \qquad \qquad \qquad \qquad \qquad \qquad \qquad \qquad \qquad \qquad \qquad \qquad \qquad \qquad \qquad \qquad \qquad \qquad \qquad \qquad \qquad \qquad \qquad \qquad \qquad \qquad \qquad \qquad \qquad \qquad \qquad \qquad \qquad \\ 0 \longrightarrow \mathrm{N} _ {1} ^ {\prime} \stackrel {{u _ {1}}} {{\longrightarrow}} \mathrm{N} _ {1} \stackrel {{v _ {1}}} {{\longrightarrow}} \mathrm{N} _ {1} ^ {\prime \prime} \longrightarrow 0 \end{array}\tag{$(\mathcal{E}_1)$}
\]

一个具有正合行的左 A-模同态的交换图。k-模的图

\[
\begin{array}{c} \operatorname{Tor} ^ {\mathbf {A}} (\mathbf {M}, \mathbf {N} ^ {\prime \prime}) \xrightarrow {\hat {c} (\mathbf {M} , \mathcal {E})} \operatorname{Tor} ^ {\mathbf {A}} (\mathbf {M}, \mathbf {N} ^ {\prime}) \\ \operatorname{Tor} ^ {\mathbf {A}} (f, g ^ {\prime \prime}) \Big \downarrow \\ \operatorname{Tor} ^ {\mathbf {A}} (\mathbf {M} _ {1}, \mathbf {N} _ {1} ^ {\prime \prime}) \xrightarrow {\hat {c} (\mathbf {M} _ {1} , \mathcal {E} _ {1})} \operatorname{Tor} ^ {\mathbf {A}} (\mathbf {M} _ {1}, \mathbf {N} _ {1} ^ {\prime}) \end{array}
\]

是交换的。

这由将 X，第 31 页，命题 2 应用于该交换图得出

\[
\begin{array}{l} 0 \longrightarrow \mathrm{L(M)} \otimes_ {\mathrm{A}} \mathrm{N} ^ {\prime} \xrightarrow {1 \otimes u} \mathrm{L(M)} \otimes_ {\mathrm{A}} \mathrm{N} \xrightarrow {1 \otimes v} \mathrm{L(M)} \otimes_ {\mathrm{A}} \mathrm{N} ^ {\prime \prime} \longrightarrow 0 \\ \quad \mathrm{L(f)} \otimes g ^ {\prime} \Biggl \downarrow \quad \mathrm{L(f)} \otimes g \Biggl \downarrow \quad \mathrm{L(f)} \otimes g ^ {\prime \prime} \Biggl \downarrow \\ 0 \longrightarrow \mathrm{L(M} _ {1}) \otimes_ {\mathrm{A}} \mathrm{N} _ {1} ^ {\prime} \xrightarrow {1 \otimes u _ {1}} \mathrm{L(M} _ {1}) \otimes_ {\mathrm{A}} \mathrm{N} _ {1} \xrightarrow {1 \otimes v _ {1}} \mathrm{L(M} _ {1}) \otimes_ {\mathrm{A}} \mathrm{N} _ {1} ^ {\prime \prime} \longrightarrow 0. \end{array}
\]

类似地，若 N 是一个左 A-模且

\[
0 \to \mathbf {M} ^ {\prime} \xrightarrow {r} \mathbf {M} \xrightarrow {s} \mathbf {M} ^ {\prime \prime} \to 0\tag{F}
\]

一个右 A-模的正合序列，定义连接同态

\[
\partial (\mathcal {F}, \mathrm{N}): \operatorname{Tor} ^ {\mathrm{A}} \left(\mathrm{M} ^ {\prime \prime}, \mathrm{N}\right)\rightarrow \operatorname{Tor} ^ {\mathrm{A}} \left(\mathrm{M} ^ {\prime}, \mathrm{N}\right)
\]

\[
\partial_ {n} (\mathcal {F}, \mathrm{N}): \operatorname{Tor} _ {n} ^ {\mathrm{A}} \left(\mathrm{M} ^ {\prime \prime}, \mathrm{N}\right)\rightarrow \operatorname{Tor} _ {n - 1} ^ {\mathrm{A}} \left(\mathrm{M} ^ {\prime}, \mathrm{N}\right)
\]

由 \(\partial(\mathcal{F},\mathrm{N})=\overline{\psi}_{\mathrm{N}}(\mathrm{M}^{\prime})^{-1}\circ\partial(\mathcal{F}^{\mathrm{N}})\circ\overline{\psi}_{\mathrm{N}}(\mathrm{M}^{\prime\prime})\) 定义，其中 \(\partial(\mathcal{F}^{\mathrm{N}})\) 是正合序列的连接同态

\[
(\mathcal {F} ^ {N})
\]

\[
0 \rightarrow \mathrm{M} ^ {\prime} \otimes_ {\mathrm{A}} \mathrm{L} (\mathrm{N}) \rightarrow \mathrm{M} \otimes_ {\mathrm{A}} \mathrm{L} (\mathrm{N}) \rightarrow \mathrm{M} ^ {\prime \prime} \otimes_ {\mathrm{A}} \mathrm{L} (\mathrm{N}) \rightarrow 0
\]

由 \(\mathcal{F}\) 导出的，并且我们有：

定理 1 bis. --- k-模同态的向左无界序列

\[
\begin{array}{l}\rightarrow \operatorname{Tor} _ {n} ^ {\mathsf {A}} (\mathbf {M} ^ {\prime}, \mathbf {N}) \xrightarrow {\operatorname{Tor} _ {n} ^ {\mathsf {A}} (r , 1)} \operatorname{Tor} _ {n} ^ {\mathsf {A}} (\mathbf {M}, \mathbf {N}) \xrightarrow {\operatorname{Tor} _ {n} ^ {\mathsf {A}} (s , 1)} \operatorname{Tor} _ {n} ^ {\mathsf {A}} (\mathbf {M} ^ {\prime \prime}, \mathbf {N}) \xrightarrow {\partial_ {n} (\mathcal {F} , \mathbf {N})} \operatorname{Tor} _ {n - 1} ^ {\mathsf {A}} (\mathbf {M} ^ {\prime}, \mathbf {N})\\\dots \rightarrow \operatorname{Tor} _ {1} ^ {\mathsf {A}} (\mathbf {M} ^ {\prime \prime}, \mathbf {N}) \xrightarrow {\partial_ {1} (\mathcal {F} , \mathbf {N})} \mathbf {M} ^ {\prime} \otimes_ {\mathbf {A}} \mathbf {N} \xrightarrow {r \otimes 1} \mathbf {M} \otimes_ {\mathbf {A}} \mathbf {N} \xrightarrow {s \otimes 1} \mathbf {M} ^ {\prime \prime} \otimes_ {\mathbf {A}} \mathbf {N} \longrightarrow 0\end{array}
\]

是正合的。

我们留给读者来陈述并证明与定理1的推论及命题9相类似的性质。此外：

命题10。——设( \(F^{\circ}\) ) 表示左 A-模的正合序列

\[
0 \rightarrow \mathrm{M} ^ {\prime \circ} \xrightarrow {r} \mathrm{M} ^ {\circ} \xrightarrow {s} \mathrm{M} ^ {\prime \prime \circ} \rightarrow 0.
\]

图表

\[
\begin{array}{c} \operatorname{Tor} ^ {\mathbf {A}} (\mathbf {M} ^ {\prime \prime}, \mathbf {N}) \xrightarrow {\partial (\mathcal {F} , \mathbf {N})} \operatorname{Tor} ^ {\mathbf {A}} (\mathbf {M} ^ {\prime}, \mathbf {N}) \\ \sigma_ {\mathbf {M} ^ {\prime}, \mathbf {N}} \Biggl \downarrow \qquad \qquad \qquad \qquad \qquad \qquad \qquad \qquad \qquad \sigma_ {\mathbf {M} ^ {\prime}, \mathbf {N}} \Biggl \downarrow \\ \operatorname{Tor} ^ {\mathbf {A} ^ {\circ}} (\mathbf {N} ^ {\circ}, \mathbf {M} ^ {\prime \prime^ {\circ}}) \xrightarrow {\partial (\mathbf {N} ^ {\circ} , \mathcal {F} ^ {\circ})} \operatorname{Tor} ^ {\mathbf {A} ^ {\circ}} (\mathbf {N} ^ {\circ}, \mathbf {M} ^ {\prime \circ}) \end{array}
\]

是交换的。

事实上，这由将 X，第 31 页，命题 2 应用于交换图表得出

\[
\begin{array}{l} 0 \longrightarrow \mathrm{M} ^ {\prime} \otimes_ {\mathrm{A}} \mathrm{L} (\mathrm{N}) \xrightarrow {r \otimes 1} \mathrm{M} \otimes_ {\mathrm{A}} \mathrm{L} (\mathrm{N}) \xrightarrow {s \otimes 1} \mathrm{M} ^ {\prime \prime} \otimes_ {\mathrm{A}} \mathrm{L} (\mathrm{N}) \longrightarrow 0 \\ \sigma (\mathrm{M} ^ {\prime}, \mathrm{L} (\mathrm{N})) \Biggl \downarrow \quad \sigma (\mathrm{M}, \mathrm{L} (\mathrm{N})) \Biggl \downarrow \quad \sigma (\mathrm{M} ^ {\prime \prime}, \mathrm{L} (\mathrm{N})) \Biggl \downarrow \\ 0 \longrightarrow \mathrm{L} (\mathrm{N} ^ {\circ}) \otimes_ {\mathrm{A} ^ {\circ}} \mathrm{M} ^ {\prime \circ} \xrightarrow {1 \otimes r} \mathrm{L} (\mathrm{N} ^ {\circ}) \otimes_ {\mathrm{A} ^ {\circ}} \mathrm{M} ^ {\circ} \xrightarrow {1 \otimes s} \mathrm{L} (\mathrm{N} ^ {\circ}) \otimes_ {\mathrm{A} ^ {\circ}} \mathrm{M} ^ {\prime \prime \circ} \longrightarrow 0. \end{array}
\]

我们将在后面看到其他的交换关系（参见 X，第 131 页，推论 1）。

